Este capítulo describe la estrategia técnica y operativa para el despliegue del sistema DOMU en un entorno productivo, asegurando su disponibilidad, seguridad y correcta adopción por parte de la comunidad.

\section{Arquitectura de Despliegue}

Para la puesta en marcha se ha diseñado una arquitectura basada en contenedores Docker, lo que garantiza la portabilidad del sistema y facilita el escalado horizontal. La Figura~\ref{fig:despliegue} ilustra los componentes del despliegue.

\begin{figure}[H]
\centering
\begin{tabular}{|p{0.92\linewidth}|}
\hline
\textbf{Diagrama de Despliegue --- DOMU} \\[4pt]
\texttt{[Cliente Web]} $\longrightarrow$ \texttt{CDN (Vercel/Netlify)} $\longrightarrow$ \texttt{SPA React} \\[2pt]
\texttt{[Cliente Móvil]} $\longrightarrow$ \texttt{PWA (navegador móvil, instalable en inicio)} \\[2pt]
\texttt{[Conserjería]} $\longrightarrow$ \texttt{Navegador + Pistola QR} \\[6pt]
\hfill $\downarrow$ \texttt{HTTPS / WSS} $\downarrow$ \hfill \\[4pt]
\texttt{[Servidor de Aplicaciones]} \\
\quad \texttt{Docker: eclipse-temurin:21-alpine} \\
\quad \texttt{Javalin 6 (API REST + WebSocket)} \\
\quad \texttt{Puerto 7070} \\[4pt]
\hfill $\downarrow$ \texttt{JDBC} $\downarrow$ \hfill \\[4pt]
\texttt{[Base de Datos]} \texttt{MySQL 8.0 (AWS RDS / Cloud SQL)} \\[4pt]
\texttt{[Almacenamiento]} \texttt{Box (archivos, comprobantes, imágenes)} \\
\hline
\end{tabular}
\caption{Arquitectura de despliegue del sistema DOMU.}
\label{fig:despliegue}
\end{figure}

Los componentes principales del despliegue son los siguientes:

\begin{itemize}
    \item \textbf{Servidor de Aplicaciones:} El backend Java se empaqueta en una imagen Docker basada en \texttt{eclipse-temurin:21-alpine}, optimizada para bajo consumo de recursos. Se despliega en un servicio de orquestación como AWS ECS o en un VPS con Docker Compose. La imagen expone el puerto~7070 y se configura mediante variables de entorno para los parámetros de base de datos, JWT, correo electrónico y almacenamiento.

    \item \textbf{Base de Datos:} Se utiliza una instancia gestionada de MySQL~8.0 en la nube (AWS RDS o Google Cloud SQL) para garantizar copias de seguridad automáticas, alta disponibilidad y mantenimiento delegado al servicio administrado. La conexión desde el servidor de aplicaciones se realiza mediante JDBC con HikariCP como pool de conexiones.

    \item \textbf{Frontend Web:} La aplicación React se construye como una SPA (\textit{Single Page Application}) y se distribuye a través de una red de entrega de contenido (CDN) como Vercel o Netlify. El proceso de \textit{build} genera archivos estáticos optimizados que se sirven directamente desde la CDN, reduciendo la latencia para los usuarios finales.

    \item \textbf{Aplicación Móvil (PWA):} La experiencia móvil se implementa como una \textit{Progressive Web App} (PWA). Los usuarios acceden desde el navegador de su dispositivo móvil y pueden instalar la aplicación en la pantalla de inicio para una experiencia similar a una app nativa, sin necesidad de distribución por tiendas de aplicaciones. La PWA se sirve desde la misma CDN que el frontend web, simplificando el despliegue y la actualización.

    \item \textbf{Almacenamiento de Archivos:} Los archivos estáticos (evidencias, comprobantes de pago, fotos de perfil, imágenes del marketplace) se alojan en Box, configurado con políticas de acceso privado y URLs firmadas temporalmente para la descarga segura.
\end{itemize}

\subsection{Configuración Docker}

El servidor de aplicaciones se despliega mediante un archivo \texttt{docker-compose.yml} que define los servicios, redes y volúmenes necesarios. La configuración incluye:

\begin{itemize}
    \item Servicio \texttt{domu-api}: imagen Docker del backend con variables de entorno para la configuración de base de datos, JWT, SMTP y Box.
    \item Servicio \texttt{mysql}: instancia de MySQL~8.0 con volumen persistente para los datos (utilizado en desarrollo; en producción se reemplaza por una instancia gestionada).
    \item Red interna \texttt{domu-network}: aísla la comunicación entre los servicios del sistema.
\end{itemize}

\subsection{Gestión de Migraciones}

Al iniciar el contenedor del backend, Flyway ejecuta automáticamente las migraciones de base de datos pendientes (35~migraciones versionadas, V001 a V035). Este mecanismo asegura que el esquema de la base de datos esté siempre sincronizado con la versión del código desplegado, eliminando la necesidad de scripts manuales de actualización.

\section{Plan de Capacitación}

La adopción del sistema requiere nivelar las competencias digitales de los distintos usuarios. Se diseñó un plan de capacitación diferenciado por perfil de usuario.

\subsection{Administración y Comité}
Se realiza una sesión de inducción de 2~horas enfocada en la configuración inicial del edificio: carga de la nómina de residentes, configuración de alícuotas, definición de áreas comunes, creación de períodos de gastos comunes y generación de estados de cuenta. Para el Comité, la capacitación considera explícitamente los flujos de gobernanza comunitaria implementados en la plataforma, manteniendo consistencia con los casos de uso y diseño funcional definidos en capítulos previos.
\begin{itemize}
    \item \textbf{Generación de votaciones:} creación de procesos de votación, definición de opciones, cierre y publicación de resultados desde el módulo de votaciones.
    \item \textbf{Carga de actas de asamblea:} publicación de documentos en biblioteca con la etiqueta obligatoria \texttt{actas}, siguiendo la trazabilidad documental del sistema.
    \item \textbf{Consulta de estados de cuenta:} revisión de estados de cuenta en modo fiscalización para seguimiento financiero comunitario.
\end{itemize}
Se entrega un manual de gestión en formato digital con procedimientos separados para Administración y Comité, incluyendo prácticas operativas y criterios de uso por rol.

\subsection{Personal de Conserjería}
Capacitación práctica en terreno (30~minutos) centrada exclusivamente en el uso del módulo de Control de Accesos y la operación de la pistola lectora de QR. Se prioriza la rapidez operativa para no entorpecer el flujo de ingreso de personas. La capacitación incluye:
\begin{itemize}
    \item Inicio de sesión y navegación básica.
    \item Registro de visitas mediante escaneo QR de cédula de identidad.
    \item Verificación de pre-autorizaciones de visitas.
    \item Registro manual de visitas sin cédula (visitantes extranjeros o cédulas deterioradas).
    \item Recepción y registro de encomiendas.
\end{itemize}

\subsection{Residentes}
Se distribuye un video tutorial de 3~minutos vía WhatsApp y correo electrónico, explicando cómo descargar la aplicación, registrarse, recuperar la contraseña, pre-autorizar una visita, registrar una reserva de área común y consultar el estado de gastos comunes. Adicionalmente, se habilita una sección de ayuda dentro de la aplicación con preguntas frecuentes.

\section{Estrategia de Lanzamiento}

Se recomienda una estrategia de \textbf{marcha blanca} de dos semanas, diseñada para minimizar el riesgo de la transición y permitir una adopción gradual del sistema.

\begin{enumerate}
    \item \textbf{Semana~1 (Paralelo):} El sistema DOMU funciona simultáneamente con el libro de visitas manual y los procesos vigentes. No se bloquean accesos ni se cobran multas asociadas al sistema digital. El objetivo es poblar la base de datos con residentes, unidades habitacionales y contactos frecuentes, y habituar al personal de conserjería a la operación del escáner QR.

    \item \textbf{Semana~2 (Transición):} Se elimina el registro en papel para visitas y encomiendas. Se habilita el módulo de reservas de áreas comunes. Los residentes comienzan a recibir notificaciones y a interactuar con la plataforma. En paralelo, Administración y Comité ejecutan una validación operativa guiada de los flujos de votaciones, carga de actas y consulta de estados de cuenta. Se monitorean incidencias y se resuelven en tiempo real.

    \item \textbf{Lanzamiento Oficial:} Se habilita el módulo financiero para el siguiente ciclo de facturación mensual. Se genera el primer estado de cuenta a través del sistema, se activan los flujos de pago y queda operativo el uso regular de los flujos de gobernanza por parte del Comité.
\end{enumerate}

\section{Respaldo y Recuperación}

Para garantizar la integridad de la información, se implementa la siguiente estrategia de respaldo:

\begin{itemize}
    \item \textbf{Respaldo de base de datos:} Se configuran rutinas automáticas (\textit{cron jobs}) que ejecutan un respaldo completo diario (\texttt{mysqldump}) de la base de datos MySQL. Los archivos de respaldo se comprimen y almacenan en un \textit{bucket} de almacenamiento en la nube (Google Drive o AWS~S3) con retención de 30~días.

    \item \textbf{Respaldo de archivos:} Los documentos almacenados en Box cuentan con redundancia y versionado nativo del servicio, garantizando la recuperación ante eliminaciones accidentales.

    \item \textbf{Objetivo de punto de recuperación (RPO):} 24~horas. En caso de fallo, se puede restaurar la base de datos al estado del último respaldo diario.

    \item \textbf{Objetivo de tiempo de recuperación (RTO):} 2~horas. El procedimiento de recuperación consiste en: (1)~restaurar el respaldo en una nueva instancia de MySQL, (2)~verificar la integridad de los datos, y (3)~reapuntar el servidor de aplicaciones a la nueva instancia.
\end{itemize}

\section{Monitoreo y Mantenimiento}

\begin{itemize}
    \item \textbf{Logs de aplicación:} El servidor Javalin registra todas las solicitudes HTTP y eventos del sistema en la salida estándar del contenedor Docker. Los logs se recopilan mediante el driver de logging de Docker y pueden integrarse con servicios de monitoreo como CloudWatch o Datadog.

    \item \textbf{Métricas de salud:} Se expone un endpoint \texttt{/health} que verifica la conectividad con la base de datos y el estado del servidor. Este endpoint es utilizado por el orquestador de contenedores para detectar y reiniciar instancias no saludables.

    \item \textbf{Monitoreo funcional de implantación:} Durante la marcha blanca se supervisan además los flujos críticos por perfil, incluyendo las operaciones de Administración y Comité sobre votaciones, biblioteca documental (etiqueta \texttt{actas}) y consulta de estados de cuenta.

    \item \textbf{Mantenimiento correctivo:} Las correcciones de errores (\textit{hotfixes}) se despliegan mediante la construcción de una nueva imagen Docker y la actualización del servicio en el orquestador, proceso que no requiere tiempo de inactividad gracias al despliegue \textit{rolling}.

    \item \textbf{Mantenimiento evolutivo:} Las nuevas funcionalidades se desarrollan en ramas independientes, se prueban en un entorno de staging y se despliegan en producción siguiendo el mismo proceso de actualización de imágenes Docker.
\end{itemize}

\section{Plan de Rollback}

En caso de que el despliegue presente fallos críticos que afecten la operación del sistema, se ejecuta el siguiente procedimiento de reversión:

\begin{enumerate}
    \item Revertir la imagen Docker del servidor de aplicaciones a la versión anterior mediante el tag de la imagen previamente etiquetada.
    \item Si la nueva versión incluyó migraciones de base de datos incompatibles, restaurar el respaldo diario más reciente de la base de datos.
    \item Verificar el correcto funcionamiento del sistema con la versión anterior, incluyendo validación de los flujos críticos de Administración y Comité.
    \item Comunicar a los usuarios afectados sobre la interrupción temporal del servicio.
\end{enumerate}

\section{Limitaciones del Despliegue}

Se identifican las siguientes limitaciones en la estrategia de despliegue de la versión~1.0:

\begin{itemize}
    \item \textbf{Ausencia de pipeline CI/CD:} El proceso de construcción y despliegue de imágenes Docker se realiza de forma manual. La implementación de un pipeline automatizado mediante GitHub Actions se contempla como trabajo futuro.
    \item \textbf{Limitaciones de la PWA:} Al ser una Progressive Web App, algunas funcionalidades nativas (como notificaciones push en iOS pre-16.4 y acceso a cierto hardware) no están disponibles. Sin embargo, la cobertura de navegadores modernos es suficiente para el público objetivo del sistema.
    \item \textbf{Instancia única:} En la versión~1.0, el servidor de aplicaciones opera como instancia única sin balanceo de carga. Para escenarios de alta concurrencia, se requiere configurar múltiples réplicas del servicio.
\end{itemize}
